% GENERATED FILE. Edit canonical lessons, then run npm run generate:books.
% canonical-source-hash: fnv1a64:cb9a1d3b3b6f3f65
% canonical-lessons: RU-C425-otplyvat, RU-C425-prichalivat, RU-C425-kachka, RU-C425-ukachivat, RU-C425-yakor

\chapter{Seasickness}
\label{ch:ru-seasickness}
\hlchaptermodality{425}

\begin{chapteropening}
\textbf{By the end of this chapter:} \emph{I can talk about a sea trip: setting sail, mooring, the rolling of the ship, seasickness and the anchor.}

Complete the last lesson of chapter 425: I can talk about a sea trip: setting sail, mooring, the rolling of the ship, seasickness and the anchor.
\end{chapteropening}

\section[\texorpdfstring{\ru{отплывать} (otplyvát)}{отплывать (otplyvát)}]{\ru{отплывать} (otplyvát) — to sail off, to set sail}

\label{lesson:RU-C425-otplyvat}



\begin{quote}
Before the new one: say the Russian for a mooring, then the Russian for a landing stage.
\end{quote}

\subsection*{You'll want to know: \ru{отплывать}}
\textbf{\ru{отплывать}} (\emph{otplyvát}) — \textquotedblleft{}to sail off, to set sail\textquotedblright{}.

\begin{quote}
\ru{Корабль} \ru{отплывает} \ru{ровно} \ru{в} \ru{полдень}.
\end{quote}

\begin{quote}
\emph{Korábl otplyváyet róvno v pólden.}
\end{quote}

\begin{quote}
The ship sets sail at exactly noon.
\end{quote}

Imperfective; perfective \emph{\ru{отплыть}} (\emph{\ru{отплыву}, \ru{отплывёшь}}).

\begin{cousinweb}
\emph{\ru{отплывать}} is built from \emph{\ru{от}-} 'away' and \emph{-\ru{плывать}}, from \emph{\ru{плыть}} 'to swim, to sail', which is inherited from Proto-Slavic \emph{plyti} (EtymDB).
\end{cousinweb}

\begin{grammarlens}[title={What you've built}]
One more word for telling, explaining and describing at B1.
\end{grammarlens}

\subsection*{Your turn}
\begin{itemize}
\raggedright
  \item \emph{Say it:} \emph{otplyvát}
  \item \emph{Say it:} \emph{otplyvát}, once more
  \item \emph{Say it:} a sentence of your own with \emph{otplyvát}
  \item \emph{Recall:} say the Russian for a mooring, then the Russian for a landing stage, then say \emph{otplyvát} again
\end{itemize}

\begin{tcolorbox}[breakable,colback=teal!4,colframe=teal!35!black,title={Before you move on}]
What does \ru{отплывать} mean? (\textquotedblleft{}To sail off, to set sail\textquotedblright{}.) Say it once more.
\end{tcolorbox}
\section[\texorpdfstring{\ru{причаливать} (prichálivat)}{причаливать (prichálivat)}]{\ru{причаливать} (prichálivat) — to moor, to dock}

\label{lesson:RU-C425-prichalivat}



\begin{quote}
Before the new one: say the Russian for a landing stage, then the Russian for to sail off.
\end{quote}

\subsection*{You'll want to know: \ru{причаливать}}
\textbf{\ru{причаливать}} (\emph{prichálivat}) — \textquotedblleft{}to moor, to dock\textquotedblright{}.

\begin{quote}
\ru{Паром} \ru{причаливает} \ru{к} \ru{пристани}.
\end{quote}

\begin{quote}
\emph{Paróm prichálivayet k prístani.}
\end{quote}

\begin{quote}
The ferry docks at the landing stage.
\end{quote}

Imperfective; perfective \emph{\ru{причалить}}. Takes \emph{\ru{к}} + dative.

\begin{cousinweb}
\emph{\ru{причаливать}} is built from \emph{\ru{при}-} 'to' and \emph{-\ru{чаливать}}, from \emph{\ru{чалить}} 'to tie up (a boat)'. EtymDB lists \emph{\ru{чалить}} without a source, so its origin is left open.
\end{cousinweb}

\begin{grammarlens}[title={What you've built}]
One more word for telling, explaining and describing at B1.
\end{grammarlens}

\subsection*{Your turn}
\begin{itemize}
\raggedright
  \item \emph{Say it:} \emph{prichálivat}
  \item \emph{Say it:} \emph{prichálivat}, once more
  \item \emph{Say it:} a sentence of your own with \emph{prichálivat}
  \item \emph{Recall:} say the Russian for a landing stage, then the Russian for to sail off, then say \emph{prichálivat} again
\end{itemize}

\begin{tcolorbox}[breakable,colback=teal!4,colframe=teal!35!black,title={Before you move on}]
What does \ru{причаливать} mean? (\textquotedblleft{}To moor, to dock\textquotedblright{}.) Say it once more.
\end{tcolorbox}
\section[\texorpdfstring{\ru{качка} (káchka)}{качка (káchka)}]{\ru{качка} (káchka) — the rolling (of a ship), the swell}

\label{lesson:RU-C425-kachka}



\begin{quote}
Before the new one: say the Russian for to sail off, then the Russian for to moor.
\end{quote}

\subsection*{You'll want to know: \ru{качка}}
\textbf{\ru{качка}} (\emph{káchka}) — \textquotedblleft{}the rolling (of a ship), the swell\textquotedblright{}.

\begin{quote}
\ru{Из}-\ru{за} \ru{качки} \ru{мама} \ru{не} \ru{могла} \ru{спать}.
\end{quote}

\begin{quote}
\emph{Iz-za káchki máma ne moglá spat.}
\end{quote}

\begin{quote}
Because of the rolling Mum could not sleep.
\end{quote}

Feminine. Built on \emph{\ru{качать}} — to rock.

\begin{cousinweb}
\emph{\ru{качка}} is built on \emph{\ru{качать}} 'to rock, to swing'; EtymDB lists both words without a deeper source.
\end{cousinweb}

\begin{grammarlens}[title={What you've built}]
One more word for telling, explaining and describing at B1.
\end{grammarlens}

\subsection*{Your turn}
\begin{itemize}
\raggedright
  \item \emph{Say it:} \emph{káchka}
  \item \emph{Say it:} \emph{káchka}, once more
  \item \emph{Say it:} a sentence of your own with \emph{káchka}
  \item \emph{Recall:} say the Russian for to sail off, then the Russian for to moor, then say \emph{káchka} again
\end{itemize}

\begin{tcolorbox}[breakable,colback=teal!4,colframe=teal!35!black,title={Before you move on}]
What does \ru{качка} mean? (\textquotedblleft{}The rolling (of a ship), the swell\textquotedblright{}.) Say it once more.
\end{tcolorbox}
\section[\texorpdfstring{\ru{укачивать} (ukáchivat)}{укачивать (ukáchivat)}]{\ru{укачивать} (ukáchivat) — to make seasick or carsick; to rock to sleep}

\label{lesson:RU-C425-ukachivat}



\begin{quote}
Before the new one: say the Russian for to moor, then the Russian for the rolling.
\end{quote}

\subsection*{You'll want to know: \ru{укачивать}}
\textbf{\ru{укачивать}} (\emph{ukáchivat}) — \textquotedblleft{}to make seasick or carsick; to rock to sleep\textquotedblright{}.

\begin{quote}
\ru{Меня} \ru{всегда} \ru{укачивает} \ru{на} \ru{море}.
\end{quote}

\begin{quote}
\emph{Menyá vsegdá ukáchivayet na móre.}
\end{quote}

\begin{quote}
I always get seasick at sea.
\end{quote}

Mostly impersonal with the accusative: \emph{\ru{меня} \ru{укачивает}} — I get motion sick. Perfective \emph{\ru{укачать}}.

\begin{cousinweb}
\emph{\ru{укачивать}} is built from \emph{\ru{у}-} and \emph{-\ru{качивать}}, from \emph{\ru{качать}} 'to rock': to rock someone until they are drowsy or sick. EtymDB does not list the verb.
\end{cousinweb}

\begin{grammarlens}[title={What you've built}]
One more word for telling, explaining and describing at B1.
\end{grammarlens}

\subsection*{Your turn}
\begin{itemize}
\raggedright
  \item \emph{Say it:} \emph{ukáchivat}
  \item \emph{Say it:} \emph{ukáchivat}, once more
  \item \emph{Say it:} a sentence of your own with \emph{ukáchivat}
  \item \emph{Recall:} say the Russian for to moor, then the Russian for the rolling, then say \emph{ukáchivat} again
\end{itemize}

\begin{tcolorbox}[breakable,colback=teal!4,colframe=teal!35!black,title={Before you move on}]
What does \ru{укачивать} mean? (\textquotedblleft{}To make seasick or carsick; to rock to sleep\textquotedblright{}.) Say it once more.
\end{tcolorbox}
\section[\texorpdfstring{\ru{якорь} (yákor)}{якорь (yákor)}]{\ru{якорь} (yákor) — an anchor}

\label{lesson:RU-C425-yakor}



\begin{quote}
Before the new one: say the Russian for the rolling, then the Russian for to make seasick or carsick.
\end{quote}

\subsection*{You'll want to know: \ru{якорь}}
\textbf{\ru{якорь}} (\emph{yákor}) — \textquotedblleft{}an anchor\textquotedblright{}.

\begin{quote}
\ru{Корабль} \ru{бросил} \ru{якорь} \ru{недалеко} \ru{от} \ru{берега}.
\end{quote}

\begin{quote}
\emph{Korábl brósil yákor nedalekó ot bérega.}
\end{quote}

\begin{quote}
The ship dropped anchor not far from the shore.
\end{quote}

Masculine; plural \emph{\ru{якоря}}. \emph{\ru{Бросить} \ru{якорь}} — to drop anchor.

\begin{cousinweb}
\emph{\ru{якорь}} comes from Old East Slavic \emph{ękorĭ} (EtymDB), an old loan that is usually traced, perhaps through Norse, to Greek \emph{ánkyra} 'an anchor'.
\end{cousinweb}

\begin{grammarlens}[title={What you've built}]
One more word for telling, explaining and describing at B1.
\end{grammarlens}

\subsection*{Your turn}
\begin{itemize}
\raggedright
  \item \emph{Say it:} \emph{yákor}
  \item \emph{Say it:} \emph{yákor}, once more
  \item \emph{Say it:} a sentence of your own with \emph{yákor}
  \item \emph{Recall:} say the Russian for the rolling, then the Russian for to make seasick or carsick, then say \emph{yákor} again
\end{itemize}

\begin{tcolorbox}[breakable,colback=teal!4,colframe=teal!35!black,title={Before you move on}]
What does \ru{якорь} mean? (\textquotedblleft{}An anchor\textquotedblright{}.) Say it once more.
\end{tcolorbox}
